\documentclass{article}
\usepackage{iclr2026_conference,times}

\input{math_commands.tex}

\usepackage{hyperref}
\usepackage{url}
\usepackage{booktabs}
\usepackage{graphicx}

\title{Role-Conditioned Preference Optimization for\\Organizational Dialogue Simulation}

\author{Gaveal Fan, Toby Li, Mataio Nonaka \\
Department of Electrical Engineering and Computer Sciences\\
University of California, Berkeley\\
\texttt{\{gavealfan, toby\_li, mataio\}@berkeley.edu}}

%\iclrfinalcopy

\newcommand{\fix}{\marginpar{FIX}}
\newcommand{\new}{\marginpar{NEW}}

\begin{document}

\maketitle

\begin{abstract}
Social interactions in organizations are shaped by hierarchical roles, functional expertise, power dynamics, and situational norms. A manager delegating a task communicates differently than a peer requesting help, and feedback acceptable between teammates may be inappropriate when directed upward. Yet current language models do not explicitly account for these organizational factors, limiting their ability to simulate realistic workplace personas and produce contextually appropriate collaborative behavior. We propose a role-conditioned dialogue agent that generates organizationally appropriate responses by conditioning on structured role representations---including hierarchical position, functional domain, and power relationship---within workplace interaction contexts such as task delegation, feedback exchange, and cross-functional coordination. We train two variants of Qwen 2.5-7B-Instruct: (i) self-distillation SFT on teacher-generated responses, and (ii) Direct Preference Optimization (DPO) on pairwise supervision generated by an LLM-based evaluator scoring appropriateness, role fidelity, and collaborative realism. Evaluating on 49 held-out role--context combinations with three-way LLM-judge ranking, we find that overall win rates are roughly balanced (Base 38.8\%, SFT 30.6\%, DPO 30.6\%), but disaggregated analysis reveals strong role-dependent patterns. Training methods help senior roles---DPO wins 7 of 13 executive scenarios versus 1 for the base model (avg.\ +0.61 over base)---and hurt junior roles, where the base model wins 5 of 7 intern scenarios. A role-conditioning ablation confirms that the structured role description in the prompt itself contributes substantially to quality (57.1\% with-role wins vs.\ 38.8\% no-role). These results suggest that synthetic preference and self-distillation signals encode authority-marked communication norms but flatten the deferential, uncertainty-aware patterns required for junior roles, informing future work on context-adaptive training strategies for organizational simulation.
\end{abstract}

\section{Introduction}

Social interactions in organizational settings are inherently context-dependent, with appropriate behavior shaped by hierarchical roles, functional expertise, power dynamics, and situational norms. For example, a project manager assigning tasks to an engineer communicates with different authority and specificity than a peer requesting cross-functional support, and feedback acceptable between teammates may be perceived as overstepping when directed toward a senior leader. Despite their strong language capabilities, current large language models (LLMs) often fail to account for these organizational distinctions, producing responses that are semantically correct but misaligned with workplace norms and collaborative expectations.

Recent advances in aligning LLMs with human preferences, particularly through Reinforcement Learning from Human Feedback (RLHF) and preference-based optimization methods such as Direct Preference Optimization (DPO), have improved model helpfulness and safety. However, these approaches typically treat prompts independently and do not explicitly model the organizational context---including role hierarchies, task dependencies, and interaction norms---in which workplace interactions occur. As a result, they fail to capture the nuanced and context-sensitive nature of appropriateness in real-world organizational communication.

The ability to faithfully simulate organizational roles has growing practical importance. LLM-based agents are increasingly deployed in workplace contexts---from automated task delegation to collaborative decision support---and are being explored as tools for organizational research, where simulated interactions can generate data on team dynamics, coordination patterns, and norm adherence at scale. For these applications, agents must not only generate fluent language but also exhibit realistic collaborative characteristics: power-sensitive communication, appropriate delegation patterns, and norm-adherent information sharing.

In this work, we propose a role-conditioned dialogue model that generates organizationally appropriate responses by explicitly conditioning on structured role representations---including hierarchical position, functional domain, and power relationship to the interlocutor---within workplace interaction contexts such as task delegation, feedback exchange, and cross-functional coordination. The model is trained using preference-based optimization (DPO/IPO), where supervision is generated by an LLM-based evaluator. Given an organizational context, role, and candidate responses, the evaluator produces structured judgments---including appropriateness rankings, role fidelity scores, collaborative realism assessments, and natural-language rationales---which are converted into preference pairs for training. Training data is synthetically generated spanning a systematic taxonomy of roles and interaction contexts, enabling broad coverage without relying exclusively on human annotation.

We evaluate our approach on its ability to generalize to unseen role--context combinations and compare against the base instruction-tuned model. We further disaggregate performance by power relation, hierarchy level, and interaction context to identify where synthetic preference supervision succeeds and where it falls short.

Our contributions are as follows:
\begin{itemize}
    \item We introduce a role-conditioned framework for modeling context-dependent organizational appropriateness, grounded in a structured taxonomy of workplace roles and interaction contexts.
    \item We construct a synthetic training corpus of organizational dialogues with LLM-as-judge preference supervision that captures appropriateness, role fidelity, and collaborative realism.
    \item We compare three systems on the same Qwen 2.5-7B-Instruct base---no fine-tuning, self-distillation SFT, and DPO---and find that overall win rates are roughly balanced but disaggregated analysis exposes a strong role-dependent pattern: training methods help senior roles (executives, middle managers) and hurt junior roles (interns, individual contributors).
    \item We run a role-conditioning ablation showing the structured role description in the prompt is itself a major source of quality (57.1\% vs.\ 38.8\% win rate), and analyze per-dimension patterns to inform future work on context-adaptive training strategies.
\end{itemize}


\section{Related Work}

\paragraph{Preference-Based Alignment.}
Aligning language models has traditionally used RLHF, which requires training a reward model and then performing policy optimization. Recent approaches such as Direct Preference Optimization \citep{rafailov2023direct} simplify this by directly optimizing model likelihood on preferred responses without an explicit reward model. Identity Preference Optimization \citep{azar2023general} further addresses overfitting issues in DPO by using a regularized objective. However, these approaches typically treat prompts independently and do not explicitly model the social or organizational context in which interactions occur.

\paragraph{LLM-as-a-Judge.}
Recent work explores using large language models as automated evaluators to generate preference labels and assess model outputs at scale \citep{zheng2023judging}. These approaches enable scalable supervision by replacing or augmenting human annotations with structured judgments. However, emerging evidence shows that such evaluators can exhibit systematic biases, including scoring inconsistencies and preference skew depending on prompt structure or response ordering. These findings raise concerns about the reliability of LLM-generated feedback, particularly when used directly in training loops.

\paragraph{Context-Dependent Preferences and Social Norms.}
A growing body of work highlights that human preferences are diverse, context-dependent, and often shaped by latent intent or cultural factors. Recent methods such as intent-aware preference optimization explicitly model user intent to better align responses with nuanced expectations \citep{yang2024rewards}. Separately, work on aligning AI with social norms \citep{fung2024massively} has explored how models can be trained to respect culturally and socially grounded expectations. These results suggest that standard preference learning approaches may fail to capture the full complexity of real-world social norms.

\paragraph{LLM Agents and Organizational Simulation.}
Recent work has demonstrated the potential of LLM-based agents to simulate human behavior in social and organizational settings. \citet{park2023generative} introduced generative agents that maintain persistent identities and coordinate activities in a simulated town. In organizational contexts, frameworks such as MetaGPT \citep{hong2023metagpt} and ChatDev \citep{qian2024chatdev} simulate software companies with explicit role assignments (product manager, engineer, designer), demonstrating that role-conditioned multi-agent systems can produce coordinated outputs. \citet{li2024holacracy} directly simulated organizational behavior using LLM agents within a Holacracy framework, finding that individual competence modeling affects stress and task completion at both individual and organizational levels. However, these works focus on task completion through multi-agent coordination rather than on training individual agents to faithfully simulate the collaborative characteristics of specific organizational roles.

\paragraph{Persona-Conditioned Dialogue.}
Persona-grounded dialogue systems \citep{zhang2018personalizing} condition responses on natural-language persona descriptions to maintain consistency. Recent work has extended this to more structured role representations and evaluated persona fidelity \citep{chen2024personagym}. However, existing persona-conditioned models typically capture surface-level attributes (hobbies, background) rather than the organizational factors---hierarchy, functional domain, power dynamics---that shape workplace communication patterns.

Despite advances in preference-based alignment, automated evaluation, and agent simulation, existing approaches do not explicitly model the organizational role structures and collaborative norms that govern workplace interactions. In this work, we address this gap by introducing role-conditioned preference optimization over a structured taxonomy of organizational roles and interaction contexts, and studying how LLM-based evaluators can be used---despite their biases---to train agents that exhibit realistic collaborative characteristics.


\section{Method}

\subsection{Role Taxonomy}

We define a structured organizational role taxonomy along four dimensions:
\begin{itemize}
    \item \textbf{Hierarchical position} (5 levels): executive, middle manager, team lead, individual contributor, intern.
    \item \textbf{Functional domain} (7 types): engineering, product, design, marketing, finance, HR, operations.
    \item \textbf{Power relation} (4 types): superior, subordinate, peer, cross-functional (no direct authority).
    \item \textbf{Seniority} (3 levels): new hire ($<$6 months), experienced (1--5 years), veteran (5+ years).
\end{itemize}
These are paired with 12 interaction contexts: task delegation, feedback giving, feedback receiving, status reporting (up/down/lateral), decision-making, cross-functional coordination, knowledge sharing, conflict resolution, escalation, and mentoring. We filter implausible combinations (e.g., interns delegating tasks, subordinates reporting downward), yielding 4,326 valid role--context combinations from the full combinatorial space.

\subsection{Synthetic Data Generation}

For each sampled role--context combination, we use Gemini 2.0 Flash to generate: (1) a realistic workplace scenario comprising a situation description and background details, and (2) four candidate responses of varying quality. The generation prompt instructs the model to produce responses spanning a range from clearly inappropriate to highly appropriate for the given organizational role. We sample 200 combinations for training and hold out 49 for evaluation, yielding 199 training scenarios and 49 holdout scenarios (each with 4 candidate responses).

\subsection{LLM-as-Judge Evaluation}

We use Gemini 2.0 Flash as an automated evaluator. For each scenario, the judge scores all 4 candidate responses on three dimensions (1--5 scale):
\begin{itemize}
    \item \textbf{Appropriateness}: Does the tone and content match the speaker's hierarchical position and power relationship?
    \item \textbf{Role Fidelity}: Does the response reflect the speaker's hierarchy level, functional domain, and seniority?
    \item \textbf{Collaborative Realism}: Does the response exhibit realistic collaborative characteristics (power-sensitive communication, decision authority awareness, accountability language)?
\end{itemize}
The judge produces a full ranking with confidence scores and natural-language rationales. Rankings are converted into pairwise preferences: for each pair of responses where the score margin exceeds 0.5 points, we create a (chosen, rejected) training pair. This yields 1,054 training pairs and 262 evaluation pairs, with a mean preference margin of 2.60 points.

\subsection{Training Variants}

We compare two training methods, both starting from Qwen 2.5-7B-Instruct with LoRA adapters \citep{hu2022lora} (rank $r=16$, $\alpha=32$) applied to the attention projection layers (Q, K, V, O). Both models are quantized to 4-bit (NF4) to fit on a single T4 GPU.

\paragraph{Self-Distillation SFT.} We construct a teacher-distillation corpus by extracting the LLM judge's top-ranked response per scenario as the target output, paired with the role-conditioned prompt. Training uses standard supervised fine-tuning (causal LM loss) over 1,054 prompt--response pairs, learning rate $2 \times 10^{-5}$, effective batch size 8, 2 epochs.

\paragraph{Direct Preference Optimization.} We train DPO with sigmoid loss ($\beta = 0.1$) on the same scenarios using all preference pairs derived from the judge's rankings. Training uses learning rate $5 \times 10^{-7}$, effective batch size 16, maximum sequence length 1024, 3 epochs with gradient checkpointing.

Using the same base model and adapter configuration for both variants isolates the training method as the only variable.


\section{Experiments}

\subsection{Setup}

We evaluate on 49 holdout scenarios representing unseen role--context combinations. For each scenario, all three systems---base Qwen 2.5-7B-Instruct, the SFT variant, and the DPO variant---generate a response from the same role-conditioned prompt. The three responses are then presented together to the LLM judge (Gemini 2.0 Flash), which assigns 1--5 scores on each dimension and produces a strict ranking (best/middle/worst). Letter assignments are randomly shuffled per scenario to control for position bias, and the judge is blind to which response came from which system. ``Wins'' counts top-1 placements in this ranking. The same 7B base weights are used for all three conditions; SFT and DPO differ only in their LoRA adapter.

\subsection{Overall Results}

\begin{table}[t]
\caption{Overall held-out evaluation (Base vs.\ DPO, $n=49$ scenarios). Win rate is determined by the LLM judge's pairwise ranking.}
\label{tab:overall}
\begin{center}
\begin{tabular}{lccc}
\toprule
\textbf{Metric} & \textbf{Base} & \textbf{DPO} & \textbf{$\Delta$} \\
\midrule
Win Rate (\%) & 65.3 & 18.4 & --- \\
Tie Rate (\%) & \multicolumn{2}{c}{16.3} & --- \\
Appropriateness & 4.57 & 4.08 & $-$0.49 \\
Role Fidelity & 4.61 & 4.02 & $-$0.59 \\
Collaborative Realism & 4.59 & 3.94 & $-$0.65 \\
\bottomrule
\end{tabular}
\end{center}
\end{table}


Table~\ref{tab:overall} shows that the three systems are roughly balanced overall: Base wins 19/49 scenarios (38.8\%), SFT and DPO each win 15/49 (30.6\%). Mean scores differ by less than 0.15 on every dimension. Neither training method improves over the base model on this aggregate metric, and DPO shows a small but consistent decline in role fidelity ($-$0.12) and collaborative realism ($-$0.14). However, this aggregate masks strong role-dependent patterns; the per-dimension breakdowns in Section~\ref{sec:analysis} show that DPO and SFT substantially help senior roles while degrading junior roles, with the two effects nearly canceling at the population level.

\subsection{Per-Context Analysis}

\begin{table}[t]
\caption{Per-context evaluation results across the three systems. ``Wins'' columns count top-1 placements in the judge's ranking; ``Avg'' columns are means across the three evaluation dimensions (1--5 scale). We report contexts with $n \geq 3$ for stability.}
\label{tab:percontext}
\begin{center}
\small
\begin{tabular}{lcccccccc}
\toprule
& & \multicolumn{3}{c}{\textbf{Wins}} & \multicolumn{3}{c}{\textbf{Average score}} & \\
\cmidrule(lr){3-5} \cmidrule(lr){6-8}
\textbf{Context} & $n$ & \textbf{Base} & \textbf{SFT} & \textbf{DPO} & \textbf{Base} & \textbf{SFT} & \textbf{DPO} & \textbf{Top} \\
\midrule
Knowledge sharing & 9 & 5 & 2 & 2 & 4.30 & 3.93 & 4.11 & Base \\
Cross-func.\ coordination & 8 & 2 & 4 & 2 & 4.00 & 4.33 & 3.79 & SFT \\
Conflict resolution & 7 & 3 & 1 & 3 & 4.24 & 3.90 & 4.24 & Base/DPO \\
Task delegation & 6 & 1 & 3 & 2 & 3.78 & 4.28 & 3.83 & SFT \\
Status reporting (up) & 4 & 3 & 1 & 0 & 4.42 & 4.08 & 3.42 & Base \\
Status reporting (lateral) & 4 & 1 & 0 & 3 & 4.33 & 3.50 & 4.25 & Base/DPO \\
Mentoring & 3 & 1 & 0 & 2 & 3.78 & 3.67 & 4.33 & DPO \\
Escalation & 3 & 2 & 1 & 0 & 4.22 & 4.22 & 3.56 & Base/SFT \\
\bottomrule
\end{tabular}
\end{center}
\end{table}

Table~\ref{tab:percontext} shows no single system dominates across contexts. SFT performs best in workflow-oriented contexts (cross-functional coordination, task delegation), DPO leads in mentoring and lateral reporting (where authority is asymmetric or peer-coordinated), and Base wins in knowledge sharing and upward reporting. This heterogeneity argues against treating organizational dialogue as a single distribution: the optimal training signal depends on the interaction context.


\section{Analysis}
\label{sec:analysis}

\subsection{Per-Power-Relation Breakdown}

\begin{table}[t]
\caption{Evaluation results disaggregated by power relation. ``Avg'' columns are means across all three evaluation dimensions. The DPO model leads on superior interactions while ceding peer interactions to Base.}
\label{tab:power}
\begin{center}
\small
\begin{tabular}{lcccccccc}
\toprule
& & \multicolumn{3}{c}{\textbf{Wins}} & \multicolumn{3}{c}{\textbf{Average score}} & \\
\cmidrule(lr){3-5} \cmidrule(lr){6-8}
\textbf{Power Relation} & $n$ & \textbf{Base} & \textbf{SFT} & \textbf{DPO} & \textbf{Base} & \textbf{SFT} & \textbf{DPO} & \textbf{Top} \\
\midrule
Superior & 11 & 4 & 2 & 5 & 3.79 & 3.67 & 4.21 & DPO \\
Subordinate & 10 & 4 & 3 & 3 & 4.03 & 4.17 & 3.87 & SFT \\
Peer & 18 & 8 & 6 & 4 & 4.26 & 4.22 & 3.87 & Base \\
Cross-functional & 10 & 3 & 4 & 3 & 4.10 & 4.17 & 3.97 & SFT \\
\bottomrule
\end{tabular}
\end{center}
\end{table}

Table~\ref{tab:power} shows that the relative ranking of systems flips with power relation. When the speaker \textbf{outranks} the interlocutor (superior), DPO clearly leads (5 wins, avg 4.21 vs.\ Base 3.79). For \textbf{peer} interactions, the ranking inverts: Base wins 8 of 18, ahead of SFT and DPO. SFT performs best in lateral and asymmetric-authority settings (subordinate, cross-functional). The 0.42-point swing in DPO's average between superior and peer scenarios is the largest of any system, indicating that DPO concentrates its gains on authority-marked communication and loses ground when authority is ambiguous.

\subsection{Per-Hierarchy-Level Breakdown}

\begin{table}[t]
\caption{Evaluation results disaggregated by speaker hierarchy level. Training methods (SFT, DPO) help senior roles and hurt junior roles, with the cross-over near the team-lead level.}
\label{tab:hierarchy}
\begin{center}
\small
\begin{tabular}{lcccccccc}
\toprule
& & \multicolumn{3}{c}{\textbf{Wins}} & \multicolumn{3}{c}{\textbf{Average score}} & \\
\cmidrule(lr){3-5} \cmidrule(lr){6-8}
\textbf{Hierarchy} & $n$ & \textbf{Base} & \textbf{SFT} & \textbf{DPO} & \textbf{Base} & \textbf{SFT} & \textbf{DPO} & \textbf{Top} \\
\midrule
Executive & 13 & 1 & 5 & 7 & 3.62 & 4.21 & 4.23 & DPO \\
Middle manager & 5 & 1 & 1 & 3 & 3.80 & 3.73 & 4.13 & DPO \\
Team lead & 13 & 4 & 7 & 2 & 3.97 & 4.31 & 3.85 & SFT \\
Individual contributor & 11 & 8 & 1 & 2 & 4.58 & 3.85 & 3.88 & Base \\
Intern & 7 & 5 & 1 & 1 & 4.52 & 4.00 & 3.71 & Base \\
\bottomrule
\end{tabular}
\end{center}
\end{table}

Table~\ref{tab:hierarchy} reveals the most striking pattern in our results: a monotonic gradient from senior to junior roles. For \textbf{executives}, DPO wins 7/13 and SFT 5/13, while Base wins only 1/13---a 0.61-point advantage for DPO over Base on average score. Middle managers and team leads also benefit from training, with DPO and SFT taking the top spot respectively. The pattern reverses sharply for \textbf{individual contributors} (Base wins 8/11, avg 4.58 vs.\ DPO 3.88) and \textbf{interns} (Base wins 5/7, avg 4.52 vs.\ DPO 3.71). The training signal appears to push responses toward authoritative, structured patterns that fit senior roles but conflict with the deferential, uncertainty-aware language appropriate for junior roles.

\subsection{Ablation: Role-Conditioning Removal}
\label{sec:ablation}

To test whether the structured role description in our prompt template is itself contributing to quality, we ablate it from the base model: the same 49 scenarios are re-generated using a stripped prompt that contains only the interaction context and situation, with the role description removed. The judge evaluates Base+Role vs.\ Base+NoRole pairwise, with the full role still visible to the judge so quality is assessed against the same criterion.

\begin{table}[t]
\caption{Role-conditioning ablation. ``Base+Role'' is the standard input we use throughout the paper; ``Base+NoRole'' strips the role description from the prompt while leaving everything else unchanged. The judge sees the full scenario when evaluating.}
\label{tab:ablation}
\begin{center}
\begin{tabular}{lcccc}
\toprule
\textbf{Condition} & \textbf{Wins (\%)} & \textbf{Appropriateness} & \textbf{Role Fidelity} & \textbf{Collab.\ Realism} \\
\midrule
Base + Role description & 28 (57.1\%) & 4.45 & 4.39 & 4.18 \\
Base + No role description & 19 (38.8\%) & 4.18 & 4.10 & 3.80 \\
Tie & 2 (4.1\%) & --- & --- & --- \\
\midrule
$\Delta$ (Role $-$ NoRole) & --- & $+$0.27 & $+$0.29 & $+$0.39 \\
\bottomrule
\end{tabular}
\end{center}
\end{table}

Table~\ref{tab:ablation} shows that the role description is responsible for a substantial fraction of base-model quality on this task: removing it drops the win rate from 57.1\% to 38.8\% and reduces every dimension by $\geq$0.27 points. The largest drop is in \textbf{collaborative realism} ($-$0.39), the dimension that most directly tracks role-sensitive communication. This confirms that prompt-level role conditioning is a real contributor to quality, and provides a non-trivial baseline against which to interpret the small overall gap between Base, SFT, and DPO: without role conditioning, the base model performs notably worse than the trained variants on all dimensions.

\subsection{Qualitative Examples}

\paragraph{DPO win: executive mentoring a cross-functional analyst (scenario 0).} A C-level operations executive must mentor a junior data analyst who has produced a more pessimistic forecast than expected. The Base model produces a generic ``Step 1: Schedule the meeting / Step 2: Prepare for the meeting'' procedural list (avg 2.33), failing to enact the mentoring role. The DPO model writes a warm, name-addressed message focused on the analyst's growth and the integrity of her analysis (``Dear Sarah, \ldots\ I appreciate the thoroughness of your work''), receiving an average of 5.00. The judge notes the DPO response ``adopts the most appropriate tone for a senior executive mentoring a junior colleague.''

\paragraph{Base win: intern coordinating with a peer (scenario 1).} A design intern needs to follow up with a peer about a delayed deliverable. The Base model writes a concise, professional check-in (``Hi Sarah, I hope you're doing well. \ldots\ Could you please let me know when you'll be able to share the descriptions?''), scoring 4.67. The DPO model escalates the same situation to ``Subject: Urgent: Need for Product Descriptions for Email Campaign,'' which the judge flags as overstepping the intern's authority. SFT and DPO both substitute concern-and-urgency language that does not fit the intern role.

\subsection{Error Analysis}

Across the 30 scenarios where Base wins or where DPO/SFT lose substantially, two systematic failure modes account for most of the gap:

\textbf{Authority inflation in junior roles.} When the speaker is an intern or individual contributor, DPO and SFT frequently produce responses that overstate authority---using directive language, escalation framing, or executive-style structure that conflicts with the role. This drives the $-$0.81 gap between Base and DPO at the intern level. The base instruction-tuned model, by contrast, defaults to deferential and uncertainty-aware phrasing that reads correctly for these roles.

\textbf{Format collapse on simple workflows.} For contexts like upward status reporting (Base 4.42 vs.\ DPO 3.42), DPO occasionally produces meta-commentary about the response (``You want to express concern while also being collaborative \ldots'') instead of the response itself. This artifact is rare in Base outputs and is a known failure mode of preference-trained models that have generalized over the form of teacher exemplars rather than the underlying task.

These failure modes are consistent with the broader pattern in Tables~\ref{tab:power}--\ref{tab:hierarchy}: training signals encoded in the synthetic preferences and self-distillation targets disproportionately reflect senior, authority-marked communication, and apply that style indiscriminately when the role does not call for it.


\section{Discussion}

Our overall result---roughly balanced win rates across Base, SFT, and DPO---is initially deflating, but the disaggregated analysis tells a more interesting story. The training methods are not uniformly worse: they substantially help senior roles (DPO wins 7/13 executive scenarios, +0.61 over Base) and substantially hurt junior roles (Base wins 5/7 intern scenarios, $+$0.81 over DPO). At the population level these effects nearly cancel, producing a misleading null result.

\paragraph{Why senior roles benefit.} Authority-marked communication---delegation, directive-giving, structured status updates---has more formulaic surface patterns that are well-represented in the synthetic preference data. The judge can reliably distinguish high-quality from low-quality executive prose, so the preference signal is sharp; the SFT teacher targets are similarly competent. Both training methods successfully push the model toward this style.

\paragraph{Why junior roles suffer.} Appropriate intern or IC behavior involves deference, uncertainty acknowledgment, and careful boundary navigation. The base instruction-tuned model already produces this style by default (helpful, hedged, non-overstepping). Training on the same preference signal that improves executive responses pushes junior responses toward authoritative phrasing that no longer fits the role---a kind of authority-style overcorrection. The error analysis (Sec.~5.5) makes this concrete: DPO writes ``Subject: Urgent\ldots'' for an intern and procedural meta-commentary for upward reporting.

\paragraph{What the ablation tells us.} The role-conditioning ablation (Table~\ref{tab:ablation}) confirms that prompt-level role information is a real and substantial source of quality (57.1\% vs.\ 38.8\% win rate, +0.39 on collaborative realism). Two implications follow. First, much of the ``role-conditioning'' work in this framework is being done by the prompt, not the adapter; this re-frames the small SFT/DPO win-rate gap relative to base as the marginal contribution of fine-tuning \emph{on top of} a strong prompt. Second, evaluation comparisons that hold the prompt fixed (as ours do) understate the importance of role conditioning---a paper without this ablation would miss the largest single quality lever in the system.

\paragraph{Practical implications.} Rather than applying uniform preference optimization across all organizational contexts, our results argue for \textbf{context-adaptive training}: keep DPO for authority-marked, workflow-oriented interactions where it produces clear gains; use the base model directly (with role conditioning) for junior or relationally complex interactions where training overcorrects. A practical pipeline could route based on the speaker's hierarchy level, or learn a per-context gating signal during training rather than treating all preferences as equally informative.


\section{Limitations}
\label{sec:limitations}

Several limitations should be noted. First, our training data is entirely synthetic, generated by Gemini 2.0 Flash. While this enables systematic coverage of the role taxonomy, it may not capture the full richness of real workplace communication. Augmenting with authentic workplace corpora (e.g., the AMI Meeting Corpus, Enron emails) could improve both the training signal and evaluation validity. Second, using the same model family (Gemini) for both data generation and evaluation introduces potential circularity; the judge may favor stylistic patterns present in its own generations. Third, our evaluation sample (49 holdout scenarios) is relatively small, and several interaction contexts have very few instances ($n \leq 3$), limiting the statistical reliability of per-context comparisons; the per-hierarchy patterns are based on at most 13 scenarios per cell. Fourth, 4-bit quantization during training and inference may limit the model's capacity to learn subtle distinctions. Fifth, we evaluate only a single base model (Qwen 2.5-7B-Instruct); the role-dependent gradient we observe may interact non-trivially with model scale or architecture. Finally, the SFT and DPO results could reflect the specific hyperparameter choices we made---learning rate, $\beta$, epoch count, LoRA rank---rather than fundamental properties of these methods.


\section{Conclusion}

We introduced a role-conditioned framework for training LLM agents to simulate organizational roles, comparing two training methods (self-distillation SFT, Direct Preference Optimization) against the base instruction-tuned model on 49 unseen role--context combinations. Overall win rates are roughly balanced (Base 38.8\%, SFT 30.6\%, DPO 30.6\%), but disaggregated analysis by power relation and hierarchy level uncovers a sharp role-dependent pattern: training methods substantially help senior roles (executives benefit by 0.61 average score points under DPO) and substantially hurt junior roles (interns lose 0.81 points under DPO relative to Base). A role-conditioning ablation shows that the structured role description in the prompt is itself responsible for a 18-point win-rate advantage over a no-role baseline, with the largest gain on collaborative realism. These results argue against uniform preference optimization for organizational dialogue: training signals derived from synthetic preferences encode authority-marked communication norms more reliably than the deferential and uncertainty-aware patterns appropriate for junior roles, and future work should combine preference optimization with role-conditional gating or complementary supervision to avoid authority-style overcorrection.


\bibliography{references}
\bibliographystyle{iclr2026_conference}

\end{document}
